\documentclass[../../main.tex]{subfiles}

\begin{document}

\section{引言}\label{section:ch02-引言}

闭有界区间上的有界函数的 Riemann 积分可通过与将定义域分划为有限子区间族相联系的函数逼近来定义. Riemann 积分推广到 Lebesgue 积分通过与将定义域分解为我们称之为 Lebesgue 可测集的有限族相联系的函数逼近来实现. 每个区间是 Lebesgue 可测的. Lebesgue 可测集的丰富性提供了一个函数, 它的积分比仅使用区间可能得到更好的上逼近与下逼近. 这就得到了在非常一般的定义域上 Lebesgue 可积的更大的一类函数以及具有更好性质的积分. 例如, 在十分一般的条件下, 我们将证明: 若一个函数序列逐点收敛到极限函数, 则该极限函数的积分是逼近它的函数列的积分的极限. 本章我们为即将研究的 Lebesgue 可测函数和 Lebesgue 积分建立基础: 该基础是可测集以及这样一个集合的 Lebesgue 测度的概念.

区间 $I$ 的长度 $\ell(I)$ 定义为 $I$ 的端点的差, 若 $I$ 是有界的; 定义为 $\infty$, 若 $I$ 是无界的. 长度是集函数的一个例子, 即将集族的每个集赋予一个扩充的实数的函数. 在长度的情形中, 定义域是所有区间的集合. 本章我们推广长度集函数到大的实数集族. 例如, 开集的 "长度" 将是组成它的可数个开区间的长度的和. 然而, 由区间和开集组成的集族对我们的目标仍然太有局限性. 我们构造称为 Lebesgue 可测集的集族和称为 Lebesgue 测度的该集族的集函数, 且将它记为 $m$. Lebesgue 可测集族是一个包含所有开集与所有闭集的 $\sigma$ 代数. 集函数 $m$ 具有以下三个性质.

\begin{description}
\item[区间的测度是它的长度] 每个非空区间 $I$ 是 Lebesgue 可测的且
\begin{align*}
m(I)=\ell(I).
\end{align*}
\item[测度是平移不变的] 若 $E$ 是 Lebesgue 可测的, 而 $y$ 是任意数, 则 $E$ 通过 $y$ 的平移 $E+y=\{x+y\mid x\in E\}$ 也是 Lebesgue 可测的, 且
\begin{align*}
m(E+y)=m(E).
\end{align*}
\item[测度在可数不交集的并是可数可加的] 若 $\{E_{k}\}_{k=1}^{\infty}$ 是 Lebesgue 可测集的可数不交族, 则
\begin{align*}
m\left(\bigcup_{k=1}^{\infty}E_{k}\right)=\sum_{k=1}^{\infty}m(E_{k}).
\end{align*}
\end{description}

不可能构造出定义在实数的所有集合上且具有以上三个性质的集函数. 事实上, 甚至不存在对实数的所有集合定义的具有以上前两个性质和有限可加性的集函数（见\cref{theorem:ch02-不可加}）. 我们通过在非常丰富的集类上构造具有上述三个性质的集函数来应对这个局限性. 该构造分为两步.

我们首先构造一个称为外测度的集函数, 将它记为 $m^{*}$. 它对任何集特别是对任何区间有定义. 区间的外测度是它的长度. 外测度是平移不变的. 然而, 外测度不是有限可加的. 但它在以下意义是可数次可加的: 若 $\{E_{k}\}_{k=1}^{\infty}$ 是任何可数集族, 互不相交或相交, 则
\begin{align*}
m^{*}\left(\bigcup_{k=1}^{\infty}E_{k}\right)\leqslant\sum_{k=1}^{\infty}m^{*}(E_{k}).
\end{align*}

构造的第二步是确定一个集是 Lebesgue 可测的含义, 证明 Lebesgue 可测集族是包含开集与闭集的 $\sigma$ 代数. 我们接着限制集函数 $m^{*}$ 到 Lebesgue 可测集族, 记为 $m$, 且证明 $m$ 是可数可加的. 我们称 $m$ 为 Lebesgue 测度.

\begin{exercise}\label{exercise:ch02-01}
在前三个习题中, 令 $m$ 为对 $\sigma$ 代数 $\mathcal{A}$ 中的所有集合定义且在 $[0,\infty]$ 取值的集函数. 假设 $m$ 在 $\mathcal{A}$ 中集合的可数不交族上是可数可加的.
\begin{enumerate}[(1)]
\item 证明: 若 $A$ 和 $B$ 是 $\mathcal{A}$ 的两个集, $A\subseteq B$, 则 $m(A)\leqslant m(B)$. 该性质称为单调性.
\item 证明: 若族 $\mathcal{A}$ 中存在一个集合 $A$ 满足 $m(A)<\infty$, 则 $m(\varnothing)=0$.
\item 令 $\{E_{k}\}_{k=1}^{\infty}$ 为 $\mathcal{A}$ 中集合的可数族, 证明 $m\left(\bigcup_{n=1}^{\infty}E_{n}\right)\leqslant\sum_{k=1}^{\infty}m(E_{k})$.
\item 定义在 $\mathbf{R}$ 的所有子集的集函数 $c$ 定义如下. 若 $E$ 有无穷多个成员, 定义 $c(E)$ 为 $\infty$; 而若 $E$ 是有限的, 定义 $c(E)$ 为 $E$ 中的元素的个数; 定义 $c(\varnothing)=0$. 证明 $c$ 是一个可数可加且平移不变的集函数. 该集函数称为计数测度.
\end{enumerate}
\end{exercise}

\end{document}